\section{The doubled Eisenstein series and its definite theta value}
\label{sec:sw-boundary}

Use the base \(B\), the CM extension \(E/F\), its involution \(\sigma\),
and its standard additive character \(\psi\) from
Section~\ref{sec:theta-match}. Fix \(m,n\geq1\), spaces \(A,Y\)
of dimensions \(m,n\), and their signature data \((r,s),(r',s')\).
The hypotheses allow either hermitian \(A\) and skew-hermitian \(Y\),
or skew-hermitian \(A\) and hermitian \(Y\). In both cases
\(\min(r,s)=0\). Write \(H=U(A)\), \(G=U(Y\oplus(-Y))\).
Its Siegel parabolic \(P\) fixes the diagonal maximal isotropic
subspace. Fix the prescribed unitary Hecke character \(\chi\) with
restriction \(\eta_{E/F}^{\,m}\) on the base ideles, a doubled
Darboux model with actions \(\omega\) of \(G(\mathbb A_F)\) and
\(S\) of \(H(\mathbb A_F)\), and the model's distinguished linear
functional \(\ell\). Choose a standard rational doubled frame \(b\).
In it the Gram matrix is
\(\left(\begin{smallmatrix}0&I\\ \epsilon I&0\end{smallmatrix}\right)\)
with \(\epsilon=\pm1\), and its first columns span the diagonal.
The associated subgroup \(K\) is unitary for the positive standard
inner products at infinity, with both the matrix and its inverse
integral in the finite frame. Choose the Iwasawa height \(a\) with
\(a(pk)=|\det t|_{\mathbb A_E}\) for the Levi action \(t\) of \(p\)
on the diagonal.

Write \(D\) for a fundamental domain of \(H(F)\) in
\(H(\mathbb A_F)\) with respect to any Haar measure \(\mu\).
Let \(\Phi\) be adelic Schwartz in the model such that
\(\operatorname{span}_{\mathbb C}\{\omega(k)\Phi:k\in K\}\)
is finite-dimensional. Put
\begin{align*}
 s_0&=(m-n)/2,\qquad f_0(g)=\ell(\omega(g)\Phi),\\
 f_z(g)&=a(g)^{z-s_0}f_0(g),\\
 I_\Phi(x)&=\mu(D)^{-1}\int_D
     \Theta_F\bigl(\omega(x)S(h)\Phi\bigr)\,d\mu(h).
\end{align*}
Here \(\Theta_F\) sums a test function over rational points in its
Darboux coordinates. For every such choice and every \(x\),
the Eisenstein series
\[
 E_\Phi(z,x)=\sum_{\gamma\in P(F)\backslash G(F)} f_z(\gamma x)
\]
has a meromorphic extension holomorphic at \(s_0\), with value
\begin{equation}
 E_\Phi(s_0,x)=
 \begin{cases} I_\Phi(x),& n<m,\\ 2I_\Phi(x),&m\leq n.\end{cases}
 \label{eq:sw-def-value}
\end{equation}
The measure in this identity is an average. Scaling \(\mu\) by a
positive constant cancels between its integral and its denominator.
The arguments below describe the continuation, an adelic Fourier
normalization, and the two terms in the boundary calculation leading
to the change in rank.

\subsection{Removing the global multiplier locally}

For any section \(f\) at parameter \(s_1\), write
\(f_z(g)=a(g)^{z-s_1}f(g)\). The assertion here allows any
nonnegative integer \(c\) with the prescribed restriction
\(\chi|_{\mathbb A_F^\times}=\eta_{E/F}^{\,c}\). Let \(\mathcal M\)
be the complex span of continuous functions of growth bounded by a
power of the matrix height \(h_2\) of \(G\). A linear functional
\(L\) on \(\mathcal M\) satisfies the growth condition if for every
\(N\geq0\) there is \(C_N\) bounding \(|L(v)|\) whenever
\(|v(g)|\leq h_2(g)^N\) for every \(g\).

Two analytic inputs enter the proof of the continuation.
First, in a right half-plane \(\Re z>t_0\), the coset series has
sum \(v(z)(x)\), with \(v(z)\in\mathcal M\) and \(L(v(z))\)
holomorphic for every such \(L\). Second, there are an entire
function \(A_0\) nonzero at some point, a function \(J(z)\) in
\(\mathcal M\) with \(L(J(z))\) entire, and \(t_1\in\mathbb R\)
such that on \(\Re z>t_1\) the series obtained by multiplying
each coset summand by \(A_0(z)\) sums to \(J(z)(x)\). These are
the half-plane and multiplier theorems in the source. Its
multiplier construction uses the rank-one case and induction on
the cardinality of the indexing set of a general linear block.

Let \(t=\max(t_0,t_1)\). Uniqueness of the sum on \(\Re z>t\)
gives \(J(z)(x)=A_0(z)v(z)(x)\) for all \(x\), including zeros
of \(A_0\). Use \(v(z)\) there and \(J(z)/A_0(z)\) elsewhere
away from the zeros. Values at remaining isolated zeros can be
assigned arbitrarily in taking a meromorphic representative.
Point evaluation at \(x\) satisfies the growth condition with
\(C_N=h_2(x)^N\), proving pointwise holomorphy in the right
half-plane. Now fix any \(z_*\). Since \(A_0\) is a nonzero
entire function, its power series at \(z_*\) has a first
nonzero coefficient. Factoring out the corresponding power gives
\(A_0(z)=(z-z_*)^k u(z)\) on an open neighbourhood \(U\) on
which \(u\) is holomorphic and nonvanishing. Set
\(J_*(z)=u(z)^{-1}J(z)\) there. For every \(x\) and every
\(z\in U\setminus\{z_*\}\),
\[
                  J_*(z)(x)=(z-z_*)^k E(z,x).
\]
For each allowed \(L\), the equality
\(L(J_*(z))=u(z)^{-1}L(J(z))\) proves holomorphy on \(U\).
The order \(k\) and the neighbourhood work simultaneously for
all \(x\) and all the functionals. This proves the weak
continuation claimed in \texttt{EisBL} from the two analytic
inputs.

\subsection{The rational Fourier pairing and its normalization}

Let \(V=\operatorname{Mat}_{m,n}(E)\) and choose the two elements
\(e_1,e_2\) of the \(F\)-basis of \(E\) in the datum. Use
\(Q=\{1,\ldots,m\}\times\{1,\ldots,n\}\times\{1,2\}\).
Write \(D_q\) for the elementary matrix with entry \(e_k\) in
position \((i,j)\), where \(q=(i,j,k)\). The coordinate isomorphism
\(C\colon V\to F^Q\) expands a matrix in the \(D_q\). For the
Gram matrix \(J_A\) define an \(F\)-bilinear pairing
\[
 B_A(Y,Z)=\operatorname{Tr}_{E/F}
              \operatorname{tr}(Y^t J_A\,\sigma Z).
\]
For any \(\varepsilon=\pm1\) the rational linear map
\(L_\varepsilon\colon F^Q\to F^Q\) given by
\((L_\varepsilon C Y)_q=\varepsilon B_A(Y,D_q)\) is bijective.
Indeed, vanishing of all its coordinates says
\[
 \operatorname{Tr}_{E/F}\bigl((Y^tJ_A)_{ji}\sigma(e_k)\bigr)=0
                         \qquad(i,j,k).
\]
The two \(\sigma(e_k)\) also span \(E\) over \(F\).
Nondegeneracy of the trace pairing for the separable extension
therefore gives \((Y^tJ_A)_{ji}=0\) for every \(j,i\).
Invertibility of \(J_A\) gives \(Y=0\). The source and target are
the same finite-dimensional vector space, proving bijectivity.

Extend \(C\) adelically. Its inverse sends \(c\) to
\(\sum_q c_q D_q\), so it identifies the diagonal rational
lattice \(F^Q\) with \(V\). Normalize additive Haar measure
\(\nu\) on \(V(\mathbb A_F)\) by rational covolume one.
The pullback of the product adelic measure on \(\mathbb A_F^Q\)
also has covolume one. Indeed, the inverse image under \(C\) of
a fundamental domain is a fundamental domain by the rational
lattice bijection, and its measure for the pullback is one.
Uniqueness of Haar measure up to scalar now identifies the two
measures. For adelic \(Z\) and rational \(Y\), the scalar
\(\varepsilon B_A(Y,Z)\) is \(\sum_q (CZ)_q(L_\varepsilon C Y)_q\)
under the diagonal embedding of the rational factors.

Consequently, for every adelic Schwartz function \(h\) in this
coordinate model,
\begin{equation}
 \sum_{Y\in V}\int_{V(\mathbb A_F)}
        \psi(\varepsilon B_A(Y,Z))h(Z)\,d\nu(Z)
       =\sum_{Y\in V}h(Y).
 \label{eq:sw-poisson}
\end{equation}
Make the change of variable \(c=CZ\) in each integral. With the
indicated sign convention the result is the adelic Fourier
transform of \(h\circ C^{-1}\) at \(L_\varepsilon C Y\).
The frequencies run over \(F^Q\) exactly once, by the preceding
bijection. Apply adelic Poisson summation on \(\mathbb A_F^Q\)
with standard character and rational covolume one, then reindex the
result by \(C^{-1}\). This proves \eqref{eq:sw-poisson} and explains
why this use of a Weyl element introduces no scalar of local origin.
The general coordinate Poisson formula is an input here.

\subsection{Splitting a boundary coefficient}

Take a lower-rank datum that omits index \(i_0\) in the \(n\)-dimensional
frame, with embedding \(j\colon G'\to G\) of the doubled group of
rank \(n-1\). Write \(N,N'\) for the additive Siegel unipotent
coordinates and \(n(\beta),n'(\beta')\) for their embeddings as
upper block unipotents. Let \(\pi\beta\) be the submatrix left
after removing the \(i_0\) row and column. Let \(T\in N(F)\) have
zero row and column \(i_0\), and set \(T'=\pi T\). For a function
\(B_0\) and a rational frequency \(T\) write
\[
 \mathcal F_T B_0(x)
 =\int_{N(F)\backslash N(\mathbb A_F)}
      B_0(n(\beta)x)\,
      \psi\bigl(-\operatorname{Tr}_{E/F}\operatorname{tr}(T\beta)\bigr)
      \,d\beta ,
\]
with normalized additive Haar measure. Here \(\operatorname{Tr}_{E/F}\)
also denotes the trace after adelic base change. In the equations
below the source defines this integral on a chosen fundamental
domain; its invariance under changes of domain follows from
rational periodicity. Let \(u_M(u)\) have upper Levi block
\(I+\sum_{j\ne i_0}u_j E_{i_0 j}\); the variable \(u\) runs
over an additive fundamental domain for \(E^{n-1}\) in
\(\mathbb A_E^{n-1}\), again of volume one.

For \(f_z=a^{z-s_1}f\) as above let \(E_n(z,x)\) be its coset
sum and, at the fixed \(x\), put
\[
 E'_{x}(z,g')=\sum_{q'\in P'(F)\backslash G'(F)}
                            f_z(j(q'g')x).
\]
The boundary identity in the release holds, for all \(x\), in a
common right half-plane depending on the model and section:
\begin{align}
 \int \mathcal F_T E_n(z,u_M(u)x)\,du
    &=\mathcal F_{T'} E'_{x}(z,1)+C_2(z,x),\label{eq:sw-cells}\\
 C_2(z,d_w(\tau)x)&=\exp\!\bigl(\tau(n-2z)\bigr)C_2(z,x).
 \label{eq:sw-dil}
\end{align}
The second equation holds for every real place \(w\), real
\(\tau\), and any model dilation there along the omitted
coordinate with that parameter.

To justify interchanging the coset sum with both integrals, take
fundamental domains with compact closures in the two additive
quotients. Such domains exist by cocompactness of the rational
lattices. The product \(C\) of the closures has compact image
under \((u,\beta)\mapsto n(\beta)u_M(u)x\). The right-half-plane
majorant for the section on that compact image supplies a summable
family \(b_q\) bounding the absolute value of its \(q\)-th summand,
uniformly there. The character factor is continuous on the compact
closure, with some bound \(M\). For the summand \(k_q\), including
its character, the product probability measure gives
\[
   \sum_q\iint |k_q(u,\beta)|\,d\beta\,du
                  \leq M\sum_q |b_q|<\infty .
\]
Continuity of each summand supplies measurability. Fubini and absolute
summability permit taking both integrals term by term. Rational
invariance of the Eisenstein series, the trace relation for \(T,T'\)
and rational invariance of the character imply independence of
the fundamental domains. Equivalently, these are integrals of
periodic functions for invariant measures on the additive quotients.

In the doubled frame partition the cosets by asking whether the
\(i_0\)-th column of the lower-left block of a representative
vanishes. For the vanishing class, the source's rational parabolic
orbit bijection gives a representative \(p_0j(\gamma')\) for each
lower-rank coset \(\gamma'\), with \(p_0\in P(F)\). The block
factorization used here asserts
\[
 j(h)n(\beta)u_M(u)x
  =p_1j(hn'(\pi\beta))x,\qquad
             \det(\operatorname{Levi}(p_1))=1 .
\]
The rational factor \(p_0\) and determinant-one parabolic \(p_1\)
leave \(f_z\) unchanged under left multiplication, by its
induction law. Expanding the trace as a double sum, every term
involving \(i_0\) vanishes, hence
\(\operatorname{tr}(T\beta)=\operatorname{tr}(T'\pi\beta)\).
Thus the corresponding summand is independent of \(u\) and
depends on \(\beta\) only through \(\pi\beta\). The \(u\)-integral
contributes one. Summing the lower-rank summands gives a continuous
rational-periodic function on \(N'(\mathbb A_F)\). Integration after
projection gives its normalized integral on the smaller quotient
(the fiber integration lemma in the source). This proves that the
first class gives precisely \(\mathcal F_{T'}E'_x(z,1)\).
The absolutely convergent sum over the complementary class defines
\(C_2(z,x)\), proving \eqref{eq:sw-cells}.

Here is the computation of the second class under dilation. Its
unfolding supplies a Weyl representative \(w_1\), coordinates
\(c\in\mathbb A_E^{n-1}\times\mathbb A_F\) with measure \(\kappa\),
and integrals of the form
\[
 \int f_z\bigl(w_1 c_M(c)j(h)x\bigr)\,d\kappa(c).
\]
The entire expression is a fixed scalar times a sum over lower
cosets of Fourier integrals of these inner integrals. The
coordinate identities used by the unfolding have three parts.
Writing \(d_\tau=d_w(\tau)\), the image \(j(h)\) commutes with
\(d_\tau\). There is \(p\in P(\mathbb A_F)\) such that
\(w_1d_\tau=pw_1\), with Levi determinant \(\Delta\) satisfying
\(|\Delta|_{\mathbb A_E}=e^{-2\tau}\) and \(\chi(\Delta)=1\).
Finally there is a measurable equivalence \(c\mapsto s_\tau(c)\)
with
\[
 c_M(c)d_\tau=d_\tau c_M(s_\tau(c)),\qquad
               (s_\tau)_*\kappa=e^{2n\tau}\kappa .
\]
These are the specific imported coordinate and Jacobian formulas
in the unfolded cell. The transformation law of a section at \(s_1\)
is \(f(pg)=\chi(\Delta)|\Delta|_{\mathbb A_E}^{s_1+n/2}f(g)\).
It follows, with positive real bases for all complex powers, that
\[
 f_z(pg)=
 e^{-2\tau(z-s_1)} e^{-2\tau(s_1+n/2)} f_z(g)
          =e^{-\tau(2z+n)}f_z(g).
\]
Insert \(d_\tau x\) in the inner integral. Move \(d_\tau\) to the
left past \(j(h)\) and \(c_M(c)\), apply the parabolic law to
\(w_1d_\tau=pw_1\), then push \(\kappa\) forward by \(s_\tau\).
The scalar is \(e^{-\tau(2z+n)}e^{2n\tau}=e^{\tau(n-2z)}\).
It is independent of the lower coset and its Fourier variable and
so passes unchanged through the remaining integral and sum. This
proves \eqref{eq:sw-dil}. The bounds, first orbit identification
and second orbit unfolding may have three different convergence
thresholds. Taking their maximum yields the single half-plane in
\eqref{eq:sw-cells}--\eqref{eq:sw-dil}.

\subsection{Recovering the value across the rank ranges}

We record how \eqref{eq:sw-cells} enters the induction. For
\(m=n\) the source proves a nondegenerate Whittaker integral
identity with scalar \(2\). For \(m>n\) the corresponding scalar
is \(1\). At a nondegenerate frequency the Eisenstein coefficient
unfolds to the integral over the full Siegel unipotent of the
section on the big Weyl cell against its Fourier character.
The theta coefficient is the orbital functional on the same
frequency. Their comparison in the source combines the finite
local integrals, unramified integrals, real-place Whittaker
integrals, global factorization and a representation-density
formula. The comparison fixes the scalar at nondegenerate
frequencies. To lift the equality to \emph{all} coefficients the
source separately uses the bounds for coefficient functionals,
a vanishing result for singular tests and the Fock-space
representation argument. These results are inputs to this chapter.

For clarity, consider passage from a single polynomial Fock lift
to a general \(K\)-finite datum. Fix a real place and write the
transformed \(\Phi\) as a finite sum of polynomial Gaussian factors
there times away-place factors. Lift the \(i\)-th summand to
\(\Phi_i\) in the original model. Injectivity of the model
transform on Schwartz functions gives
\(\Phi=\sum_{i=1}^k\Phi_i\), since the difference has transform
zero. Section formation and the theta integral are linear on
these finite sums. Suppose each \(E_{\Phi_i}\) has holomorphic
value \(cI_{\Phi_i}\) at \(s_0\), as supplied by the
polynomial-lift result. Choose a convergence threshold \(t_i\)
for the coset sum at the chosen \(x\). On
\(\Re z>\sum_i|t_i|\) every term converges absolutely, hence
\(E_\Phi=\sum_i E_{\Phi_i}\) there. The finite sum of the
continuations is holomorphic at \(s_0\) and gives
\(c\sum_i I_{\Phi_i}=cI_\Phi\). This is also valid for the empty
sum. Meromorphic uniqueness identifies this germ with the
continuation constructed above.

For \(m<n\) every Gram matrix represented by \(m\) rows has rank
at most \(m\). Full-rank theta coefficients therefore vanish.
The source's lower-rank step uses the boundary decomposition
\eqref{eq:sw-cells}, the covariance \eqref{eq:sw-dil}, the
restriction of the vacuum to the zero \(i_0\)-th column, and
the corresponding identity for singular orbital integrals.
The step also uses holomorphicity of the leading coefficient in
the real-group variable and the singular-test vanishing theorem.
Its conclusion transfers the value with scalar \(2\) from all
lower-rank vacuum vectors to all vacuum vectors of the given rank.
In the model a vacuum here is a finite sum of products of
archimedean Gaussians with finite Schwartz--Bruhat factors.
The induction is on \(k=n-m\). At \(k=0\) the equal-rank value
already proved applies to the vacua, which belong to the polynomial
Fock space and are \(K\)-finite. If \(k>0\), then \(n\geq2\); the
lower-rank existence theorem supplies the datum omitting \(i_0\).
Its \(m\), first space, base, character, Haar measure and
fundamental domain are unchanged. Its \(n'\) equals \(n-1\)
because its frame indices inject onto exactly the complement of
\(i_0\). Thus \(n'-m=k-1\). Apply induction to that datum and the
stated lower-rank step to the original datum. The vacuum generation
and differentiation theorem then gives the identity for all
\(K\)-finite vectors. This proves the scalar and rank bookkeeping
in \eqref{eq:sw-def-value}, conditional in this presentation on
the local integrals, singular-coefficient theorem and generation
results just specified.

The classical comparison is Shunsuke Yamana, \emph{The Siegel--Weil
formula for unitary groups}, Pacific Journal of Mathematics
\textbf{264} (2013), Theorem~2, p.~237,
\url{https://doi.org/10.2140/pjm.2013.264.235}.
Specialize his central division algebra to \(E\), so its degree is
one. His character on base ideles is the \(m\)-th power of the
quadratic character, the group on the doubled side is split in
the sense of admitting a complete polarization, and \(s_0=(m-n)/2\).
If \(r_0\) is the rational Witt index of the first space, his
convergence criterion includes \(r_0=0\), and his value theorem
treats \(m\leq n\) and \(m-r_0>n\). Definiteness at an
archimedean place forces \(r_0=0\), so these ranges cover the two
cases here. The source additionally transports between its Darboux
functional and the standard Schr\"odinger model and uses its own
normalized Haar average.

\subsection{Transport to the stabilizer used in the theta identity}

The rank identity above was stated for the standard compact subgroup
of a rational doubled frame. The theta integrals of
Section~\ref{sec:theta-match} are allowed an adapted stabilizer
\(K_1\). The frame-conjugacy lemma provides a parabolic element
\(p\) with Levi action \(t_p\), and standard subgroup
\(K_0=p^{-1}K_1p\). Suppose \(a_1\) is the height for \(K_1\).
Put \(d=|\det t_p|_{\mathbb A_E}>0\) and
\[
                 a_0(g)=d\,a_1(gp^{-1}).
\]
This is the height for \(K_0\). Indeed, if \(k\in K_0\), then
\(pkp^{-1}\in K_1\). For \(q\in P(\mathbb A_F)\) with Levi
action \(t_q\), write \(qkp^{-1}=(qp^{-1})(pkp^{-1})\).
The defining rule for \(a_1\) gives
\(a_1(qkp^{-1})=|\det(t_qt_p^{-1})|_{\mathbb A_E}\).
Multiplying by \(d\) gives \(|\det t_q|_{\mathbb A_E}\).
Every element admits the required decomposition by transporting
the Iwasawa decomposition for \(K_1\).

Take \(\Phi_0=\omega(p^{-1})\Phi\). If \(N\) is a
finite-dimensional \(K_1\)-invariant span containing \(\Phi\),
then \(\omega(p^{-1})N\) contains the \(K_0\)-span of
\(\Phi_0\), since for \(k\in K_0\),
\(\omega(k)\Phi_0=\omega(p^{-1})\omega(pkp^{-1})\Phi\).
Apply \eqref{eq:sw-def-value} for the standard data at \(xp\).
For every rational coset representative \(\gamma\) the section
there satisfies
\[
 a_0(\gamma xp)^{z-s_0}
  \ell\bigl(\omega(\gamma xp)\Phi_0\bigr)
 =d^{z-s_0}a_1(\gamma x)^{z-s_0}
                   \ell\bigl(\omega(\gamma x)\Phi\bigr).
\]
The factor \(d^{s_0-z}=\exp((s_0-z)\log d)\) is entire and
equals one at \(s_0\). Multiplying the continued standard
Eisenstein series by this factor thus gives the continued
adapted series with the same special value.
On the theta side the parabolic commutation lemma gives
\(\omega(p^{-1})S(h)\Phi=S(h)\omega(p^{-1})\Phi\) on Schwartz
functions. Hence
\(\omega(xp)S(h)\Phi_0=\omega(x)S(h)\Phi\), also giving exactly
the adapted theta average.

The input to that commutation lemma is injectivity of all the
translated theta functionals. It has a direct Fourier proof.
In Schr\"odinger coordinates, suppose
\(\Theta_F(\rho(x,y)\phi)=0\) for all \(x,y\in\mathbb A_F^\kappa\),
where \(\kappa\) is the finite coordinate set. For fixed \(x\)
divide out the nonzero scalar \(\psi(\langle x,y\rangle/2)\)
in the Schr\"odinger action. Then
\[
    \sum_{\xi\in F^\kappa}\phi(\xi+x)
                    \psi(\langle\xi,y\rangle)=0
                \qquad(y\in\mathbb A_F^\kappa).
\]
The coefficients are absolutely summable by the rational lattice
bound for the Schwartz translate. Every nonzero \(\xi\) gives a
nontrivial character on the additive quotient. Indeed choose
\(j\) with \(\xi_j\ne0\) and put in coordinate \(j\) an adele
\(a/\xi_j\) with \(\psi(a)\ne1\), using zero in the other
coordinates. Averaging over the quotient annihilates every such
character by translation invariance. Termwise averaging is
permitted by absolute summability. The zero coefficient is
\(\phi(x)\), so \(\phi=0\). In the doubled model the theta
expansion and the Heisenberg composition law transfer this
injectivity to translates of \(\ell\). The covariance of the two
actions then gives equal values of all these functionals on
\(\omega(p)S(h)\Phi\) and \(S(h)\omega(p)\Phi\), hence equality
of the vectors. This is the argument of \texttt{theta\_inj},
\texttt{dbl\_inj}, \texttt{dbl\_comm} and \texttt{sw\_transl}.

\paragraph{Source correspondence.}
The two top declarations are
\srcid{RgEis.Hd.eisBL} and
\srcid{SiegelWeilFormula.Fin.siegelWeilStd}. The local removal
of the multiplier is \texttt{BRoad.eisBL\_of}. The rational
pairing and Haar arguments are \texttt{rp\_nondeg},
\texttt{integral\_cov\_gen}, \texttt{int\_eq\_four} and
\texttt{TW.Twp.poisson}. The boundary calculation uses
\texttt{Act.absConv}, \texttt{C4.firstEq\_of},
\texttt{Act.secLaw\_of}, \texttt{Act.lhs\_eq\_tsum} and
\texttt{Act.ctSplit\_of\_parts}. The first-cell block
factorization, normalized projection of Haar measure, second-cell
parametrization and its Jacobian occur upstream in the proofs of
\texttt{key\_id}, \texttt{fiber}, \texttt{secUnf\_of} and
\texttt{cLawEx}. The rank assembly is
\texttt{Atop.mainC\_of}, \texttt{rank\_difference\_induction}
and \texttt{Atop.fieldOf}. The analytic and arithmetic inputs
identified above remain imported at this stage of the prose
translation.
